
%% inyecta los paquetet
% Packages

\usepackage{sectsty}
\let\oldunderbar\underbar
\usepackage{blindtext}
\usepackage{amsbsy}
\usepackage{amssymb}
\usepackage{mathrsfs}
\usepackage{amsmath} 
% Identidad UAM (Acuerdo 06/2012, numeral 4.1): Helvetica Neue en todas sus
% variantes como tipografía primaria del material impreso. Exige lualatex.
\usepackage{fontspec}
\setmainfont{HelveticaNeue}[
  BoldFont       = {HelveticaNeue-Bold},
  ItalicFont     = {HelveticaNeue-Italic},
  BoldItalicFont = {HelveticaNeue-BoldItalic}
]
\setsansfont{HelveticaNeue}[
  BoldFont       = {HelveticaNeue-Bold},
  ItalicFont     = {HelveticaNeue-Italic},
  BoldItalicFont = {HelveticaNeue-BoldItalic}
]
\setmonofont{Menlo}[Scale=MatchLowercase]
\newfontfamily\uamcondensed{HelveticaNeue-CondensedBold}
\usepackage[hmargin=3cm,vmargin=2.5cm]{geometry}
\usepackage[english,spanish,es-lcroman]{babel}
\spanishdecimal{.}
\usepackage{graphicx}
\usepackage{setspace}
\usepackage[hidelinks]{hyperref}
\usepackage{fancyhdr}
 \addtolength{\topmargin}{-11.87852pt}
\usepackage{sectsty}
\usepackage{import}
\usepackage{makecell}
\usepackage{ragged2e}
\usepackage{csquotes}
\usepackage{caption}
\usepackage{subfig}
\usepackage{color}
\usepackage{siunitx}
\usepackage{tikz}
\usepackage{pgfplots}
\pgfplotsset{compat=1.18}
  \usetikzlibrary{babel}
  \usetikzlibrary{shapes.geometric, arrows}
  \usetikzlibrary{arrows,positioning,patterns,decorations.pathmorphing}
\usepackage{circuitikz}
\usepackage{etoolbox}
\usepackage[
backend=biber,
style=ieee,
sorting=ynt
]{biblatex}
\addbibresource{bibliography/referencias.bib}
\usepackage{tabularx}
\usepackage[numbib]{tocbibind}
\usepackage{multicol}
\usepackage{pgfgantt}
\usepackage{enumitem}
\usepackage{float} 
\usepackage{wrapfig}
\usepackage{smartdiagram}
\usepackage{lipsum}
\usepackage{ifthen}
\usepackage{comment}
\usepackage{microtype}
\usepackage{placeins}
%% ruta de las imagenes
\graphicspath{ {images/} }

\makeatletter
  \renewcommand{\author}[2]{\gdef\insauthor{#2}\gdef\@author{#2}\gdef\insheadauthor{#1}\gdef\@author{#1}}
  \newcommand{\programme}[1]{\gdef\insprog{#1}\gdef\@programme{#1}}
  \renewcommand{\title}[2][1]{\gdef\institle{#2}\gdef\@title{#2}\gdef\insheadtitle{#2}\gdef\@title{\tiny#2}\gdef\insversion{\othertwodigits#1}\gdef\@title{#1}}
  \newcommand{\supervisor}[2]{\gdef\inssupervisor{#1}\gdef\@supervisor{#1}\gdef\insadscription{#2}\gdef\@supervisor{#2}}
  \newcommand{\cosupervisor}[2]{\gdef\inscosupervisor{#1}\gdef\@cosupervisor{#1}\gdef\inscoadscription{#2}\gdef\@cosupervisor{#2}}
  \newcommand{\Date}{\leavevmode\hbox{\twodigits\day/\twodigits\month/\the\year}}
  \def\twodigits#1{\ifnum#1<10 0\fi\the#1}
  \def\othertwodigits#1{\ifnum#1<10 0\fi#1}
\makeatother

% Identidad UAM: encabezados con rótulo en versalitas rojas (rojo Azcapotzalco,
% numeral 5.3) sobre el título en Helvetica Neue Bold.
\newcommand{\rotulo}[1]{{\color{uamred}\bfseries\footnotesize\addfontfeature{LetterSpace=8}\MakeUppercase{#1}}}
\makeatletter
\def\@makechapterhead#1{\thispagestyle{fancy}%
  \vspace*{10\p@}%
  {\parindent \z@ \raggedright \normalfont
    \ifnum \c@secnumdepth >\m@ne
      \if@mainmatter
        \rotulo{\@chapapp\ \thechapter}\par\nobreak
        \vskip 6\p@
      \fi
    \fi
    \interlinepenalty\@M
    {\huge\bfseries #1\par\nobreak}
    \vskip 8\p@
    {\color{uamred}\rule{2.5cm}{2pt}}\par\nobreak
    \vskip 24\p@
    \cfoot{}
	\rfoot{\thepage}
    }}
\def\@makeschapterhead#1{\thispagestyle{fancy}%
  \vspace*{10\p@}%
  {\parindent \z@ \raggedright \normalfont \interlinepenalty\@M
    {\huge\bfseries #1\par\nobreak}
    \vskip 8\p@
    {\color{uamred}\rule{2.5cm}{2pt}}\par\nobreak
    \vskip 24\p@}}
\def\@seccntformat#1{{\color{uamred}\csname the#1\endcsname}\quad}
\makeatother

\pagestyle{fancy}
\setlength{\headsep}{25pt}
	\rhead{\color{uamgray}\footnotesize\insheadauthor}
	\lhead{\begin{minipage}[b]{12cm}\color{uamgray}\footnotesize\insheadtitle\end{minipage}}
	\renewcommand{\headrulewidth}{0.4pt}
	\cfoot{}
	\lfoot{\color{uamgray}\ttfamily\scriptsize m.\insversion.\Date}
	\rfoot{\bfseries\thepage}
	\renewcommand{\footrulewidth}{0.4pt}
\makeatletter
	\renewcommand{\headrule}{{\color{hairline}\hrule\@height\headrulewidth\@width\headwidth\vskip-\headrulewidth}}
	\renewcommand{\footrule}{{\color{hairline}\vskip-\footruleskip\vskip-\footrulewidth
	  \hrule\@width\headwidth\@height\footrulewidth\vskip\footruleskip}}
\makeatother

\captionsetup{font=small,labelfont={bf,color=uamred},labelsep=period}

\makeatletter
	\patchcmd\chapter
	{\if@openright\cleardoublepage\else\clearpage\fi}
	{\par}% Inserting a \par here!
	{}{}
\makeatother

\makeatletter
	\renewcommand{\chaptermark}[1]{\markboth{ \ #1}{}}
\makeatother 

\makeatletter
	\renewcommand{\section}
	{\@startsection{section}{1}{0mm}
	{-\baselineskip}{0.5\baselineskip}
	{\Large\bfseries}} % the style
\makeatother

\makeatletter
	\renewcommand{\subsection}
	{\vskip 10\p@\@startsection{subsection}{2}{0mm}
	{-\baselineskip}{0.5\baselineskip}
	{\normalfont\bfseries}} % the style
\makeatother

\makeatletter
	\def\tableofcontents{\chapter{Contenido}\vskip1cm\@starttoc{toc}}
	\def\listoffigures{\chapter{Índice de Figuras}\@starttoc{lof}}
	\def\listoftables{\chapter{Índice de Tablas}\@starttoc{lot}}
    \def\thebibliography{\chapter{Referencias}} 	
\makeatother

\renewcommand{\appendixname}{Anexo}

\setlength{\parindent}{0pt}
\setlength{\headheight}{24pt}
%%%%%%%%%%%%%%%%%%%%%%%%%%%%%%%%%%%%%%%%%%%%%
%
%       Dedication environment
\newenvironment{dedication}
{
   \cleardoublepage
   \thispagestyle{empty}
   \vspace*{\stretch{1}}
   \hfill\begin{minipage}[t]{0.66\textwidth}
   \raggedright
}%
{
   \end{minipage}
   \vspace*{\stretch{3}}
   \clearpage
}
%%%%%%%%%%%%%%%%%%%%%%%%%%%%%%%%%%%%%%%%%%%%%
%
      %Glossary
%\makeglossary

% Helvetica Neue es más ancha que Computer Modern: holgura para evitar renglones desbordados
\setlength{\emergencystretch}{3em}
